% Generated draft: requires mathematical and source review.
\documentclass[11pt]{article}
\usepackage[margin=25mm]{geometry}
\usepackage{fontspec}
\setmainfont{Latin Modern Roman}
\usepackage{amsmath,amssymb,mathtools,unicode-math}
\setmathfont{Latin Modern Math}
\usepackage{enumitem,needspace,xurl,hyperref}
\setlength{\parindent}{0pt}
\setlength{\parskip}{5pt}
\setlength{\emergencystretch}{3em}
\newcommand{\N}{\mathbb{N}}
\newcommand{\Z}{\mathbb{Z}}
\newcommand{\Q}{\mathbb{Q}}
\newcommand{\R}{\mathbb{R}}
\newcommand{\C}{\mathbb{C}}
\newcommand{\F}{\mathbb{F}}
\newcommand{\A}{\mathbb{A}}
\newcommand{\PP}{\mathbb{P}}
\newcommand{\E}{\mathbb{E}}
\newcommand{\eps}{\varepsilon}
\newcommand{\dd}{\,\mathrm{d}}
\begin{document}
\section*{UnsolvedMath ID: 30003807}
\textbf{Source problem code:} catalogue entry
\subsection{Definitions and mathematical statement}
For the formulation below, take \(d\) and \(k\) to be positive integers. This is an explicit conventional interpretation: the supplied record does not specify their ranges or whether tensor degree \(k=0\) is intended.

Equip \(\C^d\) with its standard Hermitian inner product, conjugate-linear in the first argument. Define the symmetric tensor power by
\[
\operatorname{Sym}^k(\C^d)=\{v\in(\C^d)^{\otimes k}:U_\sigma v=v\text{ for every }\sigma\in S_k\},
\]
where \(S_k\) is the permutation group on \(k\) letters and \(U_\sigma\) permutes the tensor factors. Use the inner product inherited from the tensor product. For \(y\in\C^d\), write \(y^{\otimes k}=y\otimes\cdots\otimes y\), which belongs to \(\operatorname{Sym}^k(\C^d)\).

Let \(W:\operatorname{Sym}^k(\C^d)\to\operatorname{Sym}^k(\C^d)\) be a complex-linear Hermitian operator, meaning \(W=W^*\). Define the real-valued form
\[
p_W(y)=\bigl\langle y^{\otimes k},W y^{\otimes k}\bigr\rangle.
\]
The problem is to characterize the operators \(W\) for which
\[
\forall y\in\C^d,\qquad p_W(y)\geq 0.
\]
The quantifier includes \(y=0\), and the required inequality is non-strict. The source asks for a characterization without specifying its desired form, such as a finite criterion or a certificate.
\subsection{Short English statement}
Given a Hermitian operator on a symmetric tensor power of \(\C^d\), characterize when its associated form is nonnegative on every vector \(y\in\C^d\). The ranges of the dimension and tensor degree, and the intended type of characterization, are not explicitly specified in the source.
\subsection{Sources}
\begin{itemize}
\item Supplied local UnsolvedMath registry record, permanent ID \(30003807\), titled \emph{Positivity of Hermitian Forms on Symmetric Tensor Powers}, in the Oberwolfach Reports Open Problems collection; exact source problem code \texttt{catalogue entry}. No individual proposer is identified. The record gives proposed year 2018.
\item Supplied report citation: \emph{Interactions between Operator Space Theory and Quantum Probability with Applications to Quantum Information} (2019), Oberwolfach Reports, pp.\ 1293--1379, DOI \texttt{10.4171/owr/2018/22}. \url{https://doi.org/10.4171/owr/2018/22}. The supplied background describes Oberwolfach Reports as documenting mathematical workshops and their open-problem discussions. No precise page or named contributor for this problem is supplied.
\item Supplied literature reference: \emph{Refined Reznick Positivstellensatz}. \url{https://doi.org/10.22331/q-2023-05-12-1001}. The literature assessment dated \texttt{2026-08-21} labels the problem partially solved and states that Positivstellensatz results provide sum of squares certificates and bounds for positive forms, while no complete simple characterization is available for every Hermitian form in the extracted setting. This is supplied dated curation derived from the OWR record, not an independently verified literature assessment; the time-sensitive status and cited claims require independent verification.
\end{itemize}
\end{document}
